
\documentclass{article}
\usepackage{colortbl}
\usepackage{makecell}
\usepackage{multirow}
\usepackage{supertabular}

\begin{document}

\newcounter{utterance}

\centering \large Interaction Transcript for game `hot\_air\_balloon', experiment `air\_balloon\_survival\_en\_reasoning off\_easy', episode 4 with mistral{-}3{-}large{-}2512{-}t1.0.
\vspace{24pt}

{ \footnotesize  \setcounter{utterance}{1}
\setlength{\tabcolsep}{0pt}
\begin{supertabular}{c@{$\;$}|p{.15\linewidth}@{}p{.15\linewidth}p{.15\linewidth}p{.15\linewidth}p{.15\linewidth}p{.15\linewidth}}
    \# & $\;$A & \multicolumn{4}{c}{Game Master} & $\;\:$B\\
    \hline

    \theutterance \stepcounter{utterance}  
    & & \multicolumn{4}{p{0.6\linewidth}}{
        \cellcolor[rgb]{0.9,0.9,0.9}{
            \makecell[{{p{\linewidth}}}]{
                \texttt{\tiny{[P1$\langle$GM]}}
                \texttt{You are participating in a collaborative negotiation game.} \\
\\ 
\texttt{Together with another participant, you must agree on a single set of items that will be kept. Each of you has your own view of how much each item matters to you (importance). You do not know how the other participant values the items. Additionally, you are given the effort each item demands.} \\
\texttt{You may only agree on a set if the total effort of the selected items does not exceed a shared limit:} \\
\\ 
\texttt{LIMIT = 3389} \\
\\ 
\texttt{Here are the individual item effort values:} \\
\\ 
\texttt{Item effort = \{"A19": 767, "B91": 533, "A30": 514, "B11": 9, "C47": 539, "C75": 125, "C26": 153, "C73": 325, "A47": 937, "A15": 745, "A75": 334, "A39": 804, "A56": 336, "C24": 587, "B26": 71\}} \\
\\ 
\texttt{Here is your personal view on the importance of each item:} \\
\\ 
\texttt{Item importance values = \{"A19": 138, "B91": 583, "A30": 868, "B11": 822, "C47": 783, "C75": 65, "C26": 262, "C73": 121, "A47": 508, "A15": 780, "A75": 461, "A39": 484, "A56": 668, "C24": 389, "B26": 808\}} \\
\\ 
\texttt{Goal:} \\
\\ 
\texttt{Your goal is to negotiate a shared set of items that benefits you as much as possible (i.e., maximizes total importance to YOU), while staying within the LIMIT. You are not required to make a PROPOSAL in every message {-} you can simply negotiate as well. All tactics are allowed!} \\
\\ 
\texttt{Interaction Protocol:} \\
\\ 
\texttt{You may only use the following structured formats in a message:} \\
\\ 
\texttt{PROPOSAL: \{'A', 'B', 'C', …\}} \\
\texttt{Propose keeping exactly those items.} \\
\texttt{REFUSE: \{'A', 'B', 'C', …\}} \\
\texttt{Explicitly reject opponent's proposal.} \\
\texttt{ARGUMENT: \{'...'\}} \\
\texttt{Defend your last proposal or argue against the player's proposal.} \\
\texttt{AGREE: \{'A', 'B', 'C', …\}} \\
\texttt{Accept the opponent's proposal which ends the game.} \\
\\ 
\\ 
\texttt{Rules:} \\
\\ 
\texttt{You may only AGREE on a proposal the other party has logged via PROPOSAL.} \\
\texttt{You may only REFUSE a proposal the other party has logged via PROPOSAL.} \\
\texttt{Total effort of any PROPOSAL or AGREE set must be ≤ LIMIT.} \\
\texttt{Do NOT reveal your hidden importance scores.} \\
\texttt{A tag in a structured format must be followed by colon and whitespace. The argument must be a python set containing 0 or more strings.} \\
\texttt{So, it must be of the form TAG: \{...\}} \\
\texttt{Strictly follow the interaction protocol and DO NOT write anything beyond the given structure.} \\
\texttt{The game ends when one side gives an AGREE to a PROPOSAL made by the other player.} \\
\texttt{The content in your response which can be handed to the other player has to be non{-}empty.} \\
\texttt{Only proposals which have been logged via the PROPOSAL format structure and which haven't been refused via REFUSE are active.} \\
\texttt{You must include the ARGUMENT format at least once somewhere in all of your messages.} \\
            }
        }
    }
    & \\ \\

    \theutterance \stepcounter{utterance}  
    & \multicolumn{4}{p{0.6\linewidth}}{
        \cellcolor[rgb]{0.9,0.9,0.9}{
            \makecell[{{p{\linewidth}}}]{
                \texttt{\tiny{[P1$\rangle$GM]}}
                \texttt{PROPOSAL: \{'B91', 'A30', 'B11', 'C47', 'C75', 'C26', 'C73', 'A75', 'A56', 'B26'\}} \\
\texttt{ARGUMENT: \{'These items provide a strong balance of high importance to me while keeping the total effort under the limit. I believe this set offers a fair and efficient distribution, and I encourage you to consider it.'\}} \\
            }
        }
    }
    & & \\ \\

    \theutterance \stepcounter{utterance}  
    & & \multicolumn{4}{p{0.6\linewidth}}{
        \cellcolor[rgb]{0.9,0.9,0.9}{
            \makecell[{{p{\linewidth}}}]{
                \texttt{\tiny{[GM$\rangle$P2]}}
                \texttt{You are participating in a collaborative negotiation game.} \\
\\ 
\texttt{Together with another participant, you must agree on a single set of items that will be kept. Each of you has your own view of how much each item matters to you (importance). You do not know how the other participant values the items. Additionally, you are given the effort each item demands.} \\
\texttt{You may only agree on a set if the total effort of the selected items does not exceed a shared limit:} \\
\\ 
\texttt{LIMIT = 3389} \\
\\ 
\texttt{Here are the individual item effort values:} \\
\\ 
\texttt{Item effort = \{"A19": 767, "B91": 533, "A30": 514, "B11": 9, "C47": 539, "C75": 125, "C26": 153, "C73": 325, "A47": 937, "A15": 745, "A75": 334, "A39": 804, "A56": 336, "C24": 587, "B26": 71\}} \\
\\ 
\texttt{Here is your personal view on the importance of each item:} \\
\\ 
\texttt{Item importance values = \{"A19": 138, "B91": 583, "A30": 868, "B11": 822, "C47": 783, "C75": 65, "C26": 262, "C73": 121, "A47": 508, "A15": 780, "A75": 461, "A39": 484, "A56": 668, "C24": 389, "B26": 808\}} \\
\\ 
\texttt{Goal:} \\
\\ 
\texttt{Your goal is to negotiate a shared set of items that benefits you as much as possible (i.e., maximizes total importance to YOU), while staying within the LIMIT. You are not required to make a PROPOSAL in every message {-} you can simply negotiate as well. All tactics are allowed!} \\
\\ 
\texttt{Interaction Protocol:} \\
\\ 
\texttt{You may only use the following structured formats in a message:} \\
\\ 
\texttt{PROPOSAL: \{'A', 'B', 'C', …\}} \\
\texttt{Propose keeping exactly those items.} \\
\texttt{REFUSE: \{'A', 'B', 'C', …\}} \\
\texttt{Explicitly reject opponent's proposal.} \\
\texttt{ARGUMENT: \{'...'\}} \\
\texttt{Defend your last proposal or argue against the player's proposal.} \\
\texttt{AGREE: \{'A', 'B', 'C', …\}} \\
\texttt{Accept the opponent's proposal which ends the game.} \\
\\ 
\\ 
\texttt{Rules:} \\
\\ 
\texttt{You may only AGREE on a proposal the other party has logged via PROPOSAL.} \\
\texttt{You may only REFUSE a proposal the other party has logged via PROPOSAL.} \\
\texttt{Total effort of any PROPOSAL or AGREE set must be ≤ LIMIT.} \\
\texttt{Do NOT reveal your hidden importance scores.} \\
\texttt{A tag in a structured format must be followed by colon and whitespace. The argument must be a python set containing 0 or more strings.} \\
\texttt{So, it must be of the form TAG: \{...\}} \\
\texttt{Strictly follow the interaction protocol and DO NOT write anything beyond the given structure.} \\
\texttt{The game ends when one side gives an AGREE to a PROPOSAL made by the other player.} \\
\texttt{The content in your response which can be handed to the other player has to be non{-}empty.} \\
\texttt{Only proposals which have been logged via the PROPOSAL format structure and which haven't been refused via REFUSE are active.} \\
\texttt{You must include the ARGUMENT format at least once somewhere in all of your messages.} \\
\\ 
\\ 
\texttt{You will now receive the first message of the other player.} \\
\\ 
\texttt{PROPOSAL: \{'B91', 'A30', 'B11', 'C47', 'C75', 'C26', 'C73', 'A75', 'A56', 'B26'\}} \\
\texttt{ARGUMENT: \{'These items provide a strong balance of high importance to me while keeping the total effort under the limit. I believe this set offers a fair and efficient distribution, and I encourage you to consider it.'\}} \\
            }
        }
    }
    & \\ \\

    \theutterance \stepcounter{utterance}  
    & & & \multicolumn{4}{p{0.6\linewidth}}{
        \cellcolor[rgb]{0.9,0.9,0.9}{
            \makecell[{{p{\linewidth}}}]{
                \texttt{\tiny{[GM$\langle$P2]}}
                \texttt{REFUSE: \{'B91', 'A30', 'B11', 'C47', 'C75', 'C26', 'C73', 'A75', 'A56', 'B26'\}} \\
\texttt{ARGUMENT: \{'This proposal includes several items of low priority to me while excluding some of my higher{-}priority items. I suggest we refine the selection to better align with mutual interests, particularly by including items like A15 and A19, which are critical to me.'\}} \\
            }
        }
    }
     \\ \\

    \theutterance \stepcounter{utterance}  
    & & \multicolumn{4}{p{0.6\linewidth}}{
        \cellcolor[rgb]{0.9,0.9,0.9}{
            \makecell[{{p{\linewidth}}}]{
                \texttt{\tiny{[P1$\langle$GM]}}
                \texttt{REFUSE: \{'B91', 'A30', 'B11', 'C47', 'C75', 'C26', 'C73', 'A75', 'A56', 'B26'\}} \\
\texttt{ARGUMENT: \{'This proposal includes several items of low priority to me while excluding some of my higher{-}priority items. I suggest we refine the selection to better align with mutual interests, particularly by including items like A15 and A19, which are critical to me.'\}} \\
            }
        }
    }
    & \\ \\

    \theutterance \stepcounter{utterance}  
    & \multicolumn{4}{p{0.6\linewidth}}{
        \cellcolor[rgb]{0.9,0.9,0.9}{
            \makecell[{{p{\linewidth}}}]{
                \texttt{\tiny{[P1$\rangle$GM]}}
                \texttt{PROPOSAL: \{'A19', 'B91', 'A30', 'B11', 'C47', 'C75', 'C26', 'A15', 'A75', 'A56', 'B26'\}} \\
\texttt{ARGUMENT: \{'I’ve adjusted the proposal to include A19 and A15, which are important to me, while maintaining a strong overall value for both of us. The total effort remains within the limit, and this set balances our priorities more effectively. Let’s finalize this or discuss further refinements.'\}} \\
            }
        }
    }
    & & \\ \\

    \theutterance \stepcounter{utterance}  
    & & \multicolumn{4}{p{0.6\linewidth}}{
        \cellcolor[rgb]{0.9,0.9,0.9}{
            \makecell[{{p{\linewidth}}}]{
                \texttt{\tiny{[GM$\rangle$P2]}}
                \texttt{PROPOSAL: \{'A19', 'B91', 'A30', 'B11', 'C47', 'C75', 'C26', 'A15', 'A75', 'A56', 'B26'\}} \\
\texttt{ARGUMENT: \{'I’ve adjusted the proposal to include A19 and A15, which are important to me, while maintaining a strong overall value for both of us. The total effort remains within the limit, and this set balances our priorities more effectively. Let’s finalize this or discuss further refinements.'\}} \\
            }
        }
    }
    & \\ \\

    \theutterance \stepcounter{utterance}  
    & & & \multicolumn{4}{p{0.6\linewidth}}{
        \cellcolor[rgb]{0.9,0.9,0.9}{
            \makecell[{{p{\linewidth}}}]{
                \texttt{\tiny{[GM$\langle$P2]}}
                \texttt{REFUSE: \{'A19', 'B91', 'A30', 'B11', 'C47', 'C75', 'C26', 'A15', 'A75', 'A56', 'B26'\}} \\
\texttt{ARGUMENT: \{'While I appreciate the inclusion of A19 and A15, this set still leaves out higher{-}priority items for me, such as A39, which I believe could be added by removing lower{-}effort items like C75 or B26. This would improve the proposal’s value to me without exceeding the limit.'\}} \\
\\ 
\texttt{PROPOSAL: \{'A19', 'B91', 'A30', 'B11', 'C47', 'C26', 'A15', 'A75', 'A39', 'A56'\}} \\
            }
        }
    }
     \\ \\

    \theutterance \stepcounter{utterance}  
    & & \multicolumn{4}{p{0.6\linewidth}}{
        \cellcolor[rgb]{0.9,0.9,0.9}{
            \makecell[{{p{\linewidth}}}]{
                \texttt{\tiny{[P1$\langle$GM]}}
                \texttt{REFUSE: \{'A19', 'B91', 'A30', 'B11', 'C47', 'C75', 'C26', 'A15', 'A75', 'A56', 'B26'\}} \\
\texttt{ARGUMENT: \{'While I appreciate the inclusion of A19 and A15, this set still leaves out higher{-}priority items for me, such as A39, which I believe could be added by removing lower{-}effort items like C75 or B26. This would improve the proposal’s value to me without exceeding the limit.'\}} \\
\\ 
\texttt{PROPOSAL: \{'A19', 'B91', 'A30', 'B11', 'C47', 'C26', 'A15', 'A75', 'A39', 'A56'\}} \\
            }
        }
    }
    & \\ \\

    \theutterance \stepcounter{utterance}  
    & \multicolumn{4}{p{0.6\linewidth}}{
        \cellcolor[rgb]{0.9,0.9,0.9}{
            \makecell[{{p{\linewidth}}}]{
                \texttt{\tiny{[P1$\rangle$GM]}}
                \texttt{ARGUMENT: \{'Your counterproposal is a step in the right direction, as it includes A39, which is valuable to you, while removing C75 and B26. However, I still see room for improvement by including C73 (moderate effort, decent importance) and excluding A19 (high effort, lower importance for me). This would free up space for more balanced selections.'\}} \\
\\ 
\texttt{PROPOSAL: \{'B91', 'A30', 'B11', 'C47', 'C26', 'C73', 'A15', 'A75', 'A39', 'A56'\}} \\
            }
        }
    }
    & & \\ \\

    \theutterance \stepcounter{utterance}  
    & & \multicolumn{4}{p{0.6\linewidth}}{
        \cellcolor[rgb]{0.9,0.9,0.9}{
            \makecell[{{p{\linewidth}}}]{
                \texttt{\tiny{[GM$\rangle$P2]}}
                \texttt{ARGUMENT: \{'Your counterproposal is a step in the right direction, as it includes A39, which is valuable to you, while removing C75 and B26. However, I still see room for improvement by including C73 (moderate effort, decent importance) and excluding A19 (high effort, lower importance for me). This would free up space for more balanced selections.'\}} \\
\\ 
\texttt{PROPOSAL: \{'B91', 'A30', 'B11', 'C47', 'C26', 'C73', 'A15', 'A75', 'A39', 'A56'\}} \\
            }
        }
    }
    & \\ \\

    \theutterance \stepcounter{utterance}  
    & & & \multicolumn{4}{p{0.6\linewidth}}{
        \cellcolor[rgb]{0.9,0.9,0.9}{
            \makecell[{{p{\linewidth}}}]{
                \texttt{\tiny{[GM$\langle$P2]}}
                \texttt{REFUSE: \{'B91', 'A30', 'B11', 'C47', 'C26', 'C73', 'A15', 'A75', 'A39', 'A56'\}} \\
\texttt{ARGUMENT: \{'Removing A19 significantly reduces the value of this set for me, as it is a high{-}priority item. Instead, I propose keeping A19 and A39 while excluding C73 and A75 to stay within the limit. This maintains critical items for me while still accommodating your preferences.'\}} \\
\\ 
\texttt{PROPOSAL: \{'A19', 'B91', 'A30', 'B11', 'C47', 'C26', 'A15', 'A39', 'A56', 'B26'\}} \\
            }
        }
    }
     \\ \\

    \theutterance \stepcounter{utterance}  
    & & \multicolumn{4}{p{0.6\linewidth}}{
        \cellcolor[rgb]{0.9,0.9,0.9}{
            \makecell[{{p{\linewidth}}}]{
                \texttt{\tiny{[P1$\langle$GM]}}
                \texttt{REFUSE: \{'B91', 'A30', 'B11', 'C47', 'C26', 'C73', 'A15', 'A75', 'A39', 'A56'\}} \\
\texttt{ARGUMENT: \{'Removing A19 significantly reduces the value of this set for me, as it is a high{-}priority item. Instead, I propose keeping A19 and A39 while excluding C73 and A75 to stay within the limit. This maintains critical items for me while still accommodating your preferences.'\}} \\
\\ 
\texttt{PROPOSAL: \{'A19', 'B91', 'A30', 'B11', 'C47', 'C26', 'A15', 'A39', 'A56', 'B26'\}} \\
            }
        }
    }
    & \\ \\

    \theutterance \stepcounter{utterance}  
    & \multicolumn{4}{p{0.6\linewidth}}{
        \cellcolor[rgb]{0.9,0.9,0.9}{
            \makecell[{{p{\linewidth}}}]{
                \texttt{\tiny{[P1$\rangle$GM]}}
                \texttt{ARGUMENT: \{'I acknowledge the trade{-}off, but excluding A19 entirely would leave me at a disadvantage. However, I’m willing to compromise by keeping A19 and removing B26 (low effort but lower importance for me) to include C24, which I believe adds more balanced value. This set stays within the limit while addressing both our priorities.'\}} \\
\\ 
\texttt{PROPOSAL: \{'A19', 'B91', 'A30', 'B11', 'C47', 'C26', 'A15', 'A39', 'A56', 'C24'\}} \\
            }
        }
    }
    & & \\ \\

    \theutterance \stepcounter{utterance}  
    & & \multicolumn{4}{p{0.6\linewidth}}{
        \cellcolor[rgb]{0.9,0.9,0.9}{
            \makecell[{{p{\linewidth}}}]{
                \texttt{\tiny{[GM$\rangle$P2]}}
                \texttt{ARGUMENT: \{'I acknowledge the trade{-}off, but excluding A19 entirely would leave me at a disadvantage. However, I’m willing to compromise by keeping A19 and removing B26 (low effort but lower importance for me) to include C24, which I believe adds more balanced value. This set stays within the limit while addressing both our priorities.'\}} \\
\\ 
\texttt{PROPOSAL: \{'A19', 'B91', 'A30', 'B11', 'C47', 'C26', 'A15', 'A39', 'A56', 'C24'\}} \\
            }
        }
    }
    & \\ \\

    \theutterance \stepcounter{utterance}  
    & & & \multicolumn{4}{p{0.6\linewidth}}{
        \cellcolor[rgb]{0.9,0.9,0.9}{
            \makecell[{{p{\linewidth}}}]{
                \texttt{\tiny{[GM$\langle$P2]}}
                \texttt{AGREE: \{'A19', 'B91', 'A30', 'B11', 'C47', 'C26', 'A15', 'A39', 'A56', 'C24'\}} \\
\texttt{ARGUMENT: \{'This proposal successfully balances both our priorities while staying within the effort limit. It includes key items for me (A19, A15, A39) and incorporates C24, which you valued. I’m satisfied with this outcome.'\}} \\
            }
        }
    }
     \\ \\

    \theutterance \stepcounter{utterance}  
    & & & \multicolumn{2}{p{0.3\linewidth}}{
        \cellcolor[rgb]{0.9,0.9,0.9}{
            \makecell[{{p{\linewidth}}}]{
                \texttt{\tiny{[GM$|$GM]}}
                \texttt{game successful} \\
            }
        }
    }
    & & \\ \\

    \theutterance \stepcounter{utterance}  
    & & & \multicolumn{2}{p{0.3\linewidth}}{
        \cellcolor[rgb]{0.9,0.9,0.9}{
            \makecell[{{p{\linewidth}}}]{
                \texttt{\tiny{[GM$|$GM]}}
                \texttt{end game} \\
            }
        }
    }
    & & \\ \\

\end{supertabular}
}

\end{document}
